
\paragraph{Общие рекомендации}
\begin{itemize}
\item 
Свои программы следует оформлять в виде отдельных тестов (вместо того, чтобы модифицировать существующие).
Желательно, после того как программа заработает, сразу оставить несколько базовых \cvar{CHECK} проверок и
сделать локальный коммит, чтобы впоследствии легко распознавать и исправлять внесённые в дальнейшей работе ошибки.
К тому же наличие готовых тестов значительно облегчает рефакторинг кода.

При написании новых тестов следует переиспользовать уже написанный код, избегая копирования.
Для этого необходимо оформлять повторяющийся код в виде отдельных процедур и пользоваться механизмами наследования классов.

\item
Для запуска программ с использованием сеток большой размерности, а также для выполнения заданий, требующих
замера времени исполнения процедур, следует использовать релизную сборку проекта:
\secref{sec:debug_build} -- для сборки через cmake, \secref{sec:build_vscode} -- при работе в vscode.

\item
Рекомендации к иллюстрациям одномерных решений в Paraview 
см в~\secref{sec:paraview}. В частности, для одномерных задач см.~\secref{sec:paraview-1d}.
Для иллюстрации двумерного решения использовать изолинии (\secref{sec:paraview-isolines}) и трёхмерные поверхности (\secref{sec:paraview-2d}).

\item
Для отладки программ можно воспользоваться отладочной инфраструктурой проекта, опсанной в \secref{sec:debug}.
В случае, если решатель системы линейных уравнений не решает построенную матрицу, использовать функцию
\cvar{cfd::dbg::print} (\secref{sec:dbg-print-csr}) для отлаточной печати матрицы в консоль (размерность задачи должна быть небольшой).

\item
Для замера времени исполнения отдельных участков кода использовать класс \cvar{Tictoc}, описанный в \secref{sec:dbg-tictoc}.

\end{itemize}
